\section{Related Works} \label{sec:02-related}

\subsection{Foundations of Knowledge Transfer} \label{sec:02-related-kd}
\begin{foldable} % Knowledge distillation
  The pioneering idea in knowledge transfer between neural networks, intuitively, was to train a student to mimic the teacher's outputs~\cite{bucila2006model}.
  This concept evolved into \textit{knowledge distillation}~\cite{hinton2015distilling}, which proposed using \textit{soft} probabilistic outputs, which are more informative than \textit{hard} class labels.
  These soft labels act as essential inter-class hints, conveying the hidden relationships that define an expert's logic.
  While foundational research has expanded these mechanisms, they generally assume a transparent environment where agents have full access to both original training data and the teacher model~\cite{romero2015fitnets, nguyen2021knowledge}.
\end{foldable}

\subsection{Constraints in Data Availability} \label{sec:02-related-dfkd}
\begin{foldable} % Data-free KD
  When access to the primary data source is prohibited, researchers have explored data-efficient approaches.
  Few-shot KD methods such as~\cite{wang2020neural, nguyen2022black, vo2024improving} significantly reduce the required data volume but still rely on a small portion of the original source.
  Alternatively, other methods use proxy datasets to replace the source dataset~\cite{chen2021learning, tang2023distribution}.
  However, they often have a strong built-in bias that could be mismatched with the teacher's domain, as noted in~\cite{torralba2011unbiased}.
  Meanwhile, true data-free KD approaches often rely on dissecting the teacher's internal structure or leveraging its gradients as guidance to harvest knowledge~\cite{chen2019data, fang2021contrastive, do2022momentum, tran2024nayer}.
\end{foldable}

\subsection{Navigating Restricted Expert Interfaces} \label{sec:02-related-bbdfkd}
\begin{foldable} % ZSDB3
  Still, data-free methods collapse when the expert operates as a black-box, returning only hard labels as a top\nobreakdash-1 prediction.
  The earliest method to address BBDFKD is ZSDB3~\cite{wang2021zero}, which assumes class robustness and image fidelity based on the distance to the decision boundary.
  This method uses random pixel-wise uniform noise for initialization and gradient estimation.
  However, high-dimensional boundaries can be over-complex for distance-based approaches, so obtaining diverse image samples scales exponentially worse with high-resolution images or numerous classes, limiting their success to small-scale, simple datasets.
\end{foldable}

\begin{foldable} % IDEAL, DFHL-RS
  Subsequently, IDEAL~\cite{zhang2022ideal} and DFHL-RS~\cite{yuan2024data} utilize an extra generator network for more flexible data generation.
  Both assume that as the student improves during KD, its gradients could guide generator training.
  The methods employ in-class similarity and class-balance objectives for the generator, which is standard in data-free KD~\cite{chen2019data}.
  However, relying on the signals of an immature student is not reliable enough to train the generator.
  This setup has been observed to cause \textit{``overfitting of synthetic data''} to the student, resulting in possible mode collapse and insufficient diversity~\cite{yuan2024data}.
  Moreover, they are restricted to using hard labels instead of probabilities for KD, limiting the distillation signals.
\end{foldable}
